% 原书分章PDF第14页：18.5完整讨论。
\section{讨论}

我们的预言不等式表明，对于时间序列数据 $(O(t):t=1,\ldots)$ 及一大类时间序列模型，可以应用以损失为基础、基于交叉验证的集成学习；其中 $O(t)$ 可以是高维数据结构，而 $O(t)$ 对过去数据的依赖仅通过一个固定维度的汇总量 $Z(t)$ 体现。在线超级学习器可用于估计给定 $\bar O(t-1)$ 时 $O(t)$ 的条件概率分布的任意共同参数，只要该参数使某个损失函数的条件期望达到最小。预言不等式表明，在较弱的假设下，这个超级学习器将与预言选择器所选的估计量渐近等价。

这个在线超级学习器的另一个重要特征是，它本身是在线估计量，因此可以扩展到大型数据集。我们展示了它在独立同分布数据中的应用，由此得到一种适用于独立同分布数据的可扩展超级学习器；在计算上，它比常规超级学习器更易于处理，并且通过对数据的不同排序取平均，使交叉验证方案仿照 $V$ 折交叉验证。Benkeser 等（2017b）给出了实现并实际评价该在线超级学习器的数据分析和模拟，读者可参阅该文。当目标参数是 $P_0^n$ 的路径可微参数（通常为低维参数）时，就可以构建一个渐近正态且有效的在线靶向最小损失估计量（targeted minimum loss-based estimator，TMLE）（van der Laan 和 Rubin，2006；van der Laan，2008b；van der Laan 和 Rose，2011）。这样的 TMLE 依赖于一个良好的初始估计量。以在线超级学习器作为初始估计量，可以得到一个强有力的在线 TMLE。这为针对本章所定义的灵活／非参数时间序列模型的路径可微参数构建在线 TMLE 开辟了道路，下一章将讨论这一问题。
\clearpage
\endgroup
